% Folder 33, batch 04 — archivists' pages 61–80.
% Pass: Opus 5.5 (claude-opus-5-5), 2026-10-03 — first pass, unchecked against the pages by a human.
%
% Handwritten notes in blue then black ink on the rectos of typescript
% sheets (an English typescript on sheaves with values in a category,
% then a French typescript on root systems, both showing through). The
% even pages 62–80 are those typescripts' own faces, unrelated, carrying
% nothing in his hand: skipped. The connective prose is a fast hand and
% much of it is left \ill{}; the formulas mostly read.
\documentclass[11pt,a4paper]{article}
\input{../preamble/grothendieck}

\folder{33}
\batch{4}
\pages{61}{80}
\foldertitle{SGA 1965. MS [Manuscrit] brouillon : notes manuscrites (s.d.).}
\dating{1965-[vers 1971]}
\watermark{Édition de démonstration}

\begin{document}

\section{23. Th. de dualité globale pour un morphisme lissement compactifiable}

\page{61}
\note{le numéro « 23 » et le titre, soulignés, sont de sa main en tête de la page ; les paragraphes précédents (au moins jusqu'au n° 22) sont antérieurs au lot.}

Le morphisme $f : X \to Y$ sera dit \emph{lissement compactifiable} s'il existe
\uncertain{une} $Y$-\uncertain{immersion} \ill{} $X \to X'$, \ill{} $X'$
\uncertain{lisse} et \uncertain{projectif} sur $Y$ \ill{} \ill{} \ill{} $Y$ !).
\uncertain{Par exemple}, si $f$ est \uncertain{quasi-projectif} et $Y$ \ill{}
\ill{} \ill{}, \ill{} \ill{} \ill{} \ill{} $X$ \ill{} $\mathbb{P}^r_Y$,
\ill{} \ill{} \uncertain{lissement compactifiable}.

\uncertain{Supposons} $X$ \ill{} \uncertain{compactifiable}, et
\uncertain{choisissons} une \uncertain{immersion} $i : X \hookrightarrow X'$
\ill{} \ill{} de \uncertain{définition}.

\begin{tikzcd}
X \arrow[r, hook, "i"] \arrow[d, "f"'] & X' \arrow[dl, "{f'}"] \\
Y &
\end{tikzcd}

\uncertain{Définissons} un foncteur \struck{\ill{}}
\[
  \mathbb{R}f^! : D^+(Y) \longrightarrow D^+(X)
\]
par
\[
  \mathbb{R}f^!(G) = (\mathbb{R}i^!)(\mathbb{R}f'^!)
\]
\note{la seconde parenthèse du membre de droite est mal formée sur la page ; on lit $\mathbb{R}i^!$ suivi de $\mathbb{R}f'^!$ d'après la suite, sans certitude.}
où $\mathbb{R}f'^!$ a été défini \ill{}, et $\mathbb{R}i^!$ \uncertain{dans} \ill{}.
\uncertain{Compte tenu} du th. de dualité pour $f'$ et \ill{} \ill{} \ill{} \ill{}
\note{la phrase se poursuit en tête de la p. 63.}

\page{63}
\uncertain{isomorphisme} \uncertain{fonctoriel}
\[
  \mathrm{Hom}(F^\bullet, \mathbb{R}f^!(G^\bullet)) \simeq \mathrm{Hom}(\mathbb{R}_!f_*F^\bullet, G^\bullet)
\]
\[
  \bigl(F^\bullet \in \mathrm{Ob}\, D^b(X),\ G^\bullet \in \mathrm{Ob}\, D^+(Y)\bigr).
\]
\note{la formule est encadrée.}

\emph{Lorsque $X$ est lisse sur $Y$}, on a \ill{} \ill{} deux
\uncertain{définitions} de $\mathbb{R}f^!$ \uncertain{définies} \ill{}, et
$\mathbb{R}f^!$ défini \uncertain{directement} en termes de $f$ lisse,
\uncertain{dans} \ill{} \ill{} \ill{} \uncertain{transitivité}.

\marginal{à préciser}
[C'est essentiellement le th. de \struck{dualité locale} \add{pureté relatif}
pour $f : X \to Y$, cf.\ N.B. \ill{} \ill{} \ill{} \ill{} \ill{} \ill{}
\ill{} \ill{} \ill{} $f^!(X)$ avec la \uncertain{formule} de
\uncertain{transitivité} \ill{} \ill{} degrés …]. Ceci dit, le
\struck{\ill{}} \add{th.} de dualité \ill{} \ill{} \ill{} \ill{}
\uncertain{canonique} $f^!$ \ill{} \ill{} \ill{} \uncertain{défini}
\uncertain{directement} \ill{} \ill{} $f$ : (c'est \uncertain{essentiellement}
\ill{} \ill{} « th.\ de \uncertain{pureté relatif} \ill{} \ill{}
\uncertain{fondamentales} \uncertain{locales et globales} ») \ill{} \ill{}
\ill{} fait allusion déjà.
\note{les crochets et le double trait vertical dans la marge gauche, avec « à préciser », isolent le passage de « C'est essentiellement » à « degrés ». La lecture du passage reste très lacunaire.}

\page{65}
\uncertain{Pour} \struck{se} \uncertain{sentir} à l'aise, du point de \ill{}
\ill{} \struck{\ill{}} \uncertain{variances}, \ill{} \uncertain{dans} le
\uncertain{cadre} \ill{}, \uncertain{compte} que la description de
$\mathbb{R}f^!$ et de l'isomorphisme de dualité ne dépend pas,
\uncertain{essentiellement}, de la factorisation choisie. L'argument qui suit
est \add{dû}, en ce qui \uncertain{concerne} la dualité, à
\uncertain{Hartshorne} (dans le cas des faisceaux cohérents).
\ill{} \ill{} \ill{} deux \ill{} \ill{} \uncertain{telles que} $X'$, $X''$,
\ill{} \ill{}

\begin{tikzcd}[column sep=small]
 & X''' \arrow[ddl] \arrow[ddr] & \\
 & X \arrow[u] \arrow[dl] \arrow[dr] & \\
X' & & X''
\end{tikzcd}

\note{entre $X'''$ et $X$ il trace deux flèches verticales accolées, dont une au moins monte de $X$ vers $X'''$ ; on n'a reporté que celle-ci.}

\[
  X''' = X' \times_Y X'' ,
\]
et \uncertain{prenant} \ill{} \ill{} \ill{} \ill{}, \ill{}
\uncertain{isomorphisme} \uncertain{canonique} \ill{}
$\mathbb{R}^!_{i'} f'^{\#}$ et $\mathbb{R}^!_{i''} f''^{\#}$
\uncertain{compatible} \uncertain{avec} \struck{\ill{}} \ill{} \ill{}
\uncertain{dualité}, \ill{} \ill{} \ill{} et $\mathbb{R}^!_{i'''} f'''^{\#}$,
i.e.\ que \ill{} \ill{} \uncertain{ramène} \ill{} \ill{} \uncertain{situation}
\ill{} \ill{} \ill{} \ill{} \uncertain{lisse}

\begin{tikzcd}
 & X'' \arrow[r, "g"] \arrow[dr, "{f''}"'] & X' \arrow[d, "{f'}"] \\
X \arrow[ur, "{i''}"] \arrow[urr, "{i'}" description] \arrow[rr, "f"'] & & Y
\end{tikzcd}

\note{schéma redessiné à partir d'un croquis serré : les lettres sont sûres, la place exacte des étiquettes l'est moins.}
telle que $i' = g\, i''$.

Utilisant $\mathbb{R}^! f'' \cdot \mathbb{R}^! (f') \simeq (\mathbb{R}^! g)(\mathbb{R}^! f')$,
on est \uncertain{ramené} au cas où $X' = Y$ et $f' = \mathrm{id}_Y$.
\uncertain{Changeant} de notation, on est \uncertain{ramené} \ill{}
\ill{} \ill{} $g : X \to Y$ lisse, et $Y' \xrightarrow{\;i\;} Y$
\ill{}, \uncertain{et} \ill{} \ill{} \ill{} \struck{\ill{}} \ill{}
\ill{} de $Y'$ dans $X$, i.e.\ une \uncertain{section} $\sigma$ de
$X' = X \times_Y Y'$ sur $Y'$.

\begin{tikzcd}
X \arrow[r, hook, "j"] \arrow[d, "g"'] & X' \arrow[d, "{g'}"] \\
Y \arrow[r, hook, "i"'] & Y' \arrow[u, bend right=40, "\sigma"']
\end{tikzcd}

\note{dans le coin inférieur droit, $Y'$ est écrit sur un $X'$ biffé.}
$g$, $g'$ lisses, $i$, $j$, \struck{$\sigma$} immersions.

\page{67}
À \struck{\ill{}} démontrer un \uncertain{isom.} de \uncertain{foncteurs}.
\[
  (*) \qquad \mathbb{R}i^! \simeq \text{\struck{$\mathbb{R}g$}}\ \mathbb{R}(j\sigma)^!\, \mathbb{R}^! g .
\]
et une \uncertain{commutativité} de \uncertain{diagrammes} \ill{} \ill{}
\uncertain{concernant} les $\mathrm{Tr}_i$ et, $\mathrm{Tr}_g$, $\mathrm{Tr}_{j\sigma}$.
\[
  \mathbb{R}(j\sigma)^! \simeq \mathbb{R}\sigma^!\, \mathbb{R}j^!
\]
\uncertain{Prouvons} (*).
\[
  (j\sigma)^! \simeq \sigma^!\, j^!
\]
On \ill{} (*) comme \struck{\ill{}}
\[
  \mathbb{R}i^! \simeq \mathbb{R}\sigma^!\,(\mathbb{R}j^!\, \mathbb{R}^! g)
\]
\uncertain{déduit} \uncertain{des} \uncertain{isom.} \uncertain{suivants}
\[
  \begin{cases}
    (\mathbb{R}^! g')(\mathbb{R}i^!) \simeq (\mathbb{R}j^!)(\mathbb{R}^! g) & (a)\\
    \text{\struck{$\mathbb{R}\,\mathrm{id}_{D^+(Y)}$}}\ \mathbb{R}^!\,\mathrm{id}_{Y'} \simeq (\mathbb{R}\sigma^!)(\mathbb{R}^! g') & (b),
  \end{cases}
\]
d'où
\[
  \mathbb{R}i^! = \text{\struck{$\mathbb{R}i^!$}}\,(\mathbb{R}^!\,\mathrm{id}_{Y'})\,\mathbb{R}i^!
  \simeq \text{\struck{$\mathbb{R}i^!$}}\ \mathbb{R}\sigma^!\, \mathbb{R}^! g'\, \mathbb{R}i^!
\]
\[
  \simeq \mathbb{R}\sigma^!\,\mathbb{R}j^!\,\mathbb{R}^! g\ \text{\struck{$\mathbb{R}i^!$}}
\]
\note{au-dessus des deux derniers facteurs il écrit, entouré, « $\mathbb{R}\sigma^!$ » ; sous $\mathbb{R}^! g'\,\mathbb{R}i^!$ un mot biffé illisible, avec un renvoi « qui (a) ».}

\uncertain{en} \ill{} \ill{} \ill{} \uncertain{définition}. Il faut
\ill{} (a) et (b). \uncertain{(b)} \ill{} \ill{} \ill{} \ill{} \ill{}
\uncertain{des} « \ill{} \uncertain{locales et globales} ».
\ill{} \ill{} \ill{} de \uncertain{définition} (a) et (b) \ill{} \ill{}
\uncertain{compatibilité} \ill{} \ill{} \ill{} \ill{} \ill{}
(\uncertain{lorsque} l'un des \ill{} \ill{} \struck{\ill{}} \ill{} \ill{}
\ill{} \ill{} …)

\page{69}
d'autre part (a) s'écrit, \ill{} \ill{} \emph{shift}.
\ill{} degrés et \uncertain{tensorisation} par \struck{\ill{}}.
\ill{} $\underline{\Omega}^d_{X'/Y'}$
\[
  \Bigl(\text{N.B.}\quad \underline{\Omega}^d_{X'/Y'} \simeq \underline{\Omega}^d_{X/Y} \otimes_{\mathcal{O}_X} \mathcal{O}_{X'}\Bigr)
\]
\[
  g'^*(\mathbb{R}i^!) \simeq (\mathbb{R}j^!)\, g^* .
\]
Il \uncertain{résulte} donc \ill{} \ill{} \uncertain{du} \emph{\uncertain{Th.} de
\uncertain{changement} de base \uncertain{lisse}}.

C'est d'ailleurs l'isom.\ de \uncertain{foncteurs} \ill{}
\uncertain{intervient} \uncertain{dans} \ill{} \uncertain{définition}. Il
\ill{} à \uncertain{vérifier} la \uncertain{compatibilité} \ill{} les
\uncertain{traces} \ill{} \uncertain{dualité}, \ill{} \ill{} \ill{} \ill{}
\uncertain{difficultés}, et \ill{} \uncertain{laisse} \ill{} \ill{} …

\emph{Remarques}. a) \uncertain{Pour} \ill{} \ill{} \ill{} \ill{}, à
\ill{} \ill{} \ill{}, \ill{} \uncertain{vérifier} \ill{} \uncertain{laissons}
\ill{} \ill{} \ill{} \ill{} \uncertain{fonctorialité} $\mathbb{R}^! f'$, \ill{}
\ill{} \ill{} \ill{} $X'$ lisses $/Y$, \ill{} \ill{} \ill{} \ill{} \ill{}
\uncertain{de transitivité} \ill{} \ill{}. \ill{} la \uncertain{compatibilité}
\ill{} \ill{} \ill{} \uncertain{triviale}. \ill{} \ill{}
\uncertain{compatibilité} \ill{} \ill{} \ill{} \uncertain{plus loin} dans le
« \emph{formulaire} ».

\page{71}
b) Dans la \uncertain{définition} \ill{} \uncertain{foncteurs}
$\mathbb{R}^! f$, \ill{} \ill{} \ill{} \uncertain{utilisé} \ill{} \ill{}
que $f$. \ill{} \ill{} \uncertain{projectif} \ill{} « \uncertain{localement
projectif} » ; \struck{\ill{}} \ill{} \ill{} $X'$ \uncertain{lisse} \add{des} \ill{}
\ill{} \uncertain{relation} \ill{} $Y$ \ill{}. L'hyp.\ \ill{} \ill{}
« \uncertain{lisse compactifiable} » \ill{} \ill{} \ill{} \ill{}
\ill{} \ill{} \ill{} \ill{} $\mathbb{R}_!f$ ($\mathbb{R}^!f$) et \ill{} que
\struck{\ill{} \ill{}} \ill{} \ill{} \ill{} \ill{} \ill{} \ill{}
\uncertain{définition} \ill{} pour $f$ \uncertain{loc.\ de type fini}
\ill{} \uncertain{lisse}, \ill{} \ill{} \uncertain{comme} \ill{}
$\mathbb{R}_! f$ \ill{} \ill{} \ill{} $f$ \ill{} \ill{} \ill{} \ill{}
\uncertain{définition} \ill{} \ill{} \ill{}.

\ill{} \ill{} \uncertain{définition} \ill{} $\mathbb{R}_!f$
\ill{} \ill{}, \ill{} \ill{} \ill{} \ill{} \ill{}. \ill{} \ill{} \ill{} \ill{}
\ill{} \ill{} \ill{} \uncertain{définis}, \ill{}
\emph{\uncertain{Th}} :
\[
  \mathrm{Tr}_f : (\mathbb{R}_!f)(\mathbb{R}^!f) \longrightarrow \mathrm{id}_{D^+(Y)}
\]
\uncertain{qui} \uncertain{devrait} \uncertain{donner} \ill{} \ill{}
« \ill{} de \uncertain{dualité} » \ill{} \ill{} $f$ \ill{} \ill{} \ill{}
\uncertain{trivial} \ill{} \ill{}, \emph{\uncertain{isom.}} \ill{} \ill{}
\note{le reste de la page est biffé ou illisible ; plusieurs mots ont été remplacés par d'autres au-dessus de la ligne.}

\section{23 bis. L'homom.\ de dualité par un sous-préschéma \add{lisse} d'un préschéma lisse}

\page{73}
\note{le titre, souligné, est de sa main : « 23 », corrigé, est suivi de « bis » en exposant ; dans la seconde ligne du titre un mot est biffé (\struck{\ill{}}) et remplacé au-dessus de la ligne.}

\begin{tikzcd}
Y \arrow[rr, hook, "i"] \arrow[dr, "g"'] & & X \arrow[dl, "f"] \\
 & S &
\end{tikzcd}

\note{il note les dimensions relatives : « $(d)$ » sous $i$, « $(r-d)$ » à côté de $g$, « $(r)$ » à côté de $f$.}

\struck{Définition} \add{L'hom}
\[
  \mathbb{R}i_!\bigl(\mathbb{R}i^!(T_{X/S})\bigr) \simeq \mathbb{R}i_!\bigl(T_{Y/S}[2d]\bigr)
  \longrightarrow \mathbb{R}_!f_*(T_{X/S})
\]
\note{au-dessus : « $\mathbb{R}i^!$ \ill{} une \uncertain{immersion} \ill{} » et « $(\mathbb{R}i_!)\,\mathbb{R}i^!(F^\bullet) \to F^\bullet$ » ; une flèche descend de $\mathbb{R}i^!(T_{X/S})$ vers le mot « \uncertain{redéfinition} ». Le terme médian est réécrit sur une formule biffée, illisible.}
\[
  \mathbb{R}_!i_*\bigl(T_{Y/S}[-2d]\bigr) \longrightarrow \mathbb{R}_!f_*(T_{X/S})
\]
\[
  \mathbb{R}_!g_*(T_{Y/S})[-2d] \longrightarrow \mathbb{R}_!f_*(T_{X/S})
\]
\[
  R^{\,i-2d}_!\, g_*(T_{Y/S}) \longrightarrow R^{\,i}_!\, f_*(T_{X/S})
\]
\note{dans les deux exposants de la dernière ligne un premier signe est biffé et remplacé.}
Faisons $i = 2r$ \struck{\ill{}}, \ill{} :
\[
  R^{2(r-d)}_!\, g_*(T_{Y/S}) \longrightarrow T_{X/S}
\]
\note{le but se lit $T_{X/S}$ ; on attendrait un faisceau sur $S$. Transcrit tel quel.}

a) Cet \uncertain{homom.} et \uncertain{compatible} \uncertain{des}
\uncertain{homom.} \uncertain{trace}.

b) Plus \uncertain{généralement}, on a \ill{} un $G$ \uncertain{sur} $S$ un
isom.\ de \uncertain{foncteurs}
\[
  \text{\struck{$\mathbb{R}^! g^*(G) \simeq \mathbb{R}i^!\,\mathbb{R}f^!(G)$}}
\]
\[
  \mathbb{R}g^!(G) \simeq \mathbb{R}i^!\,\mathbb{R}f^!(G)
\]

\page{75}
qui \uncertain{généralise} le \emph{th.\ de pureté relatif}.
[et se \uncertain{déduit} d'ailleurs du th.\ de \uncertain{dualité globale} ---
cela devrait être \uncertain{explicité} dans SGAA \struck{\ill{}}]. \ill{}
\uncertain{en déduit} \ill{} \uncertain{formules}
\note{le crochet ouvert à « et se déduit » est accompagné de trois traits verticaux dans la marge gauche.}
\[
  \underline{H}^i_Y(f^*(G)) =
  \begin{cases}
    0 & i \neq 2d\\
    g^*(G) \otimes T^{-1}_{Y/X} & i = 2d
  \end{cases}
\]
Cet \uncertain{isomorphisme} \uncertain{se généralise} \uncertain{dans}
\ill{} \ill{} \ill{} un $G^\bullet$.
\[
  \mathbb{R}g^!(G^\bullet) \simeq \mathbb{R}i^!\bigl(\mathbb{R}f^!(G^\bullet)\bigr)
\]
On en \uncertain{déduit} un \uncertain{homom.} \uncertain{canonique}
\[
  \text{\struck{$\mathbb{R}_!g_*(\mathbb{R}g^!(G^\bullet)) \to \mathbb{R}f$}}
\]
\[
  \mathbb{R}_!i_*\bigl(\mathbb{R}g^!(G^\bullet)\bigr) \longrightarrow \mathbb{R}f^!(G^\bullet)
\]
d'où \ill{} \ill{} $\mathbb{R}_!f_*$, et \uncertain{utilisant} la
\uncertain{transitivité} :

\begin{tikzcd}[column sep=small, nodes={font=\scriptsize}]
\mathbb{R}_!g_*\bigl(\mathbb{R}g^!(G^\bullet)\bigr) \arrow[rr] \arrow[dr, "{\mathrm{Tr}_g(G^\bullet)}"'] & & \mathbb{R}_!f_*\,\mathbb{R}f^!(G^\bullet) \arrow[dl, "{\mathrm{Tr}_f(G^\bullet)}"] \\
 & G^\bullet &
\end{tikzcd}

\marginal{\uncertain{rend commutatif}}

\page{77}
On en \uncertain{conclut} que si $F^\bullet$ \ill{} \struck{\ill{}} \ill{}
\ill{} \ill{} $Y$, \uncertain{alors} \ill{} \uncertain{commutativité} dans

\begin{tikzcd}[column sep=small, nodes={font=\scriptsize}]
\mathrm{Ext}^i\bigl(Y; F^\bullet, \mathbb{R}g^!(G^\bullet)\bigr) \arrow[r, "\Theta_g"] \arrow[d, "{\Theta_i}"'] & \mathrm{Ext}^i\bigl(S; \mathbb{R}_!g_*(F^\bullet), G^\bullet\bigr) \arrow[d, "{\text{transitivité}}"] \\
\mathrm{Ext}^i\bigl(X; \mathbb{R}_!i_*(F^\bullet), \mathbb{R}f^!(G^\bullet)\bigr) \arrow[r, "\Theta_f"] & \mathrm{Ext}^i\bigl(S; \mathbb{R}_!f_*(\mathbb{R}_!i_*(F^\bullet)), G^\bullet\bigr)
\end{tikzcd}

\note{les deux flèches verticales portent le signe d'isomorphisme ; au-dessus de $\mathbb{R}g^!(G^\bullet)$ il écrit « $= \mathbb{R}i^!\,\mathbb{R}f^!(G^\bullet)$ » ; dans le premier $\mathrm{Ext}$, un « $X$, » est biffé devant $Y$ ; sous le coin inférieur gauche une égalité $\mathbb{R}f^!(G^\bullet) = \mathbb{R}i^!\,\mathbb{R}\ldots$ est rayée.}

i.e.\ \uncertain{compatibilité} \ill{} \ill{} de l'isom.\ de
\uncertain{dualité} $\Theta_i$ (cas \uncertain{trivial}) et de l'isom.\ de
\uncertain{transitivité} pour $f$, $i$ (cas trivial aussi) à l'isom.\ de
dualité $\Theta_g$ \ill{} « \uncertain{induit} » \ill{} $\Theta_f$.
\uncertain{Kif-kif} en \uncertain{localisant} sur la base …

\page{79}
On, ainsi \struck{des \uncertain{isomorphismes}} \add{des lois covariantes} \ill{}
\ill{} \ill{}. C'est très explicite pour $f$ lisse. \uncertain{Pour}
$f$ \uncertain{général}, \uncertain{pour} l'\uncertain{explicitation}, \ill{}
\ill{} \uncertain{servira} \ill{} \ill{} \ill{} \ill{}

2°) \emph{L'hom.\ de \uncertain{transitivité}}
\[
  \mathbb{R}^!(gf) \simeq \mathbb{R}^!f\ \mathbb{R}^!g
\]
(\uncertain{Pour mémoire}, $f$ plus \uncertain{haut}).

3°) \uncertain{Revenons} \ill{} \ill{} : 1°) \uncertain{lorsque} $X$, $Y$
\ill{} \ill{} $Z$, soit $F^\bullet \in \mathrm{Ob}\, D^+(Z)$, et
\uncertain{prenons}
\[
  G^\bullet = \mathbb{L}g^*(F^\bullet)
\]
d'où :
\ill{} $g \to$ \ill{}
\[
  G^\bullet \simeq g^!(F^\bullet)\ \text{\struck{$[-2d]$}}\otimes T_{Y/Z}\,[-2d]
\]
\note{au-dessus de la ligne, en partie barré : « \ill{} \ill{} \uncertain{rel.} $d$ ».}
\[
  f^!(G^\bullet) = \bigl[f^!\, g^!(F^\bullet)\bigr] \otimes T^{-1}_{Y/Z}\,[-2d]
\]
\note{une accolade sous $f^!\,g^!(F^\bullet)$ renvoie à la ligne suivante.}
\[
  f^!\, g^!(F^\bullet) \simeq h^!(F^\bullet) = \mathbb{L}h^*(F^\bullet) \otimes T_{X/Z}\,[2d']
\]
\note{le facteur $T_{X/Z}$ est réécrit sur un signe biffé ; en marge droite : « \struck{\ill{}} \uncertain{$d'$ dim. relative de} \ill{} ».}
\[
  f^!(G^\bullet) \simeq \mathbb{L}h^*(F^\bullet) \otimes T_{X/Y}\,[2d_{X/Y}]
\]
\note{dans le crochet final, un signe et un indice sont biffés et corrigés.}
où
\[
  T_{X/Y} \simeq T_{X/Z} \otimes T^{-1}_{Y/Z} , \qquad
  d_{X/Y} = \dim_{X/Z} - \dim_{Y/Z}
\]
$=$ dim.\ \uncertain{relative} \uncertain{absolue}.
\note{la page s'arrête ici ; la suite du 3°) n'est pas dans le lot.}

\end{document}
